\documentclass{ximera}
\title{Circles and Quadratic Equations, Part 1}


\newcommand{\pskip}{\vskip 0.1 in}

\begin{document}
\begin{abstract}
Questions about circles and Quadratic Equations.
\end{abstract}
\maketitle


\pskip

\section{Quadratic Functions}

\begin{question} \label{Qdffrp3r}
Let 
\[
   f(x) = (x-6)^2 + 5 .
\]
\begin{enumerate}

\item Evaluate $f(-1)$.

\item Solve the equation 
\[
             f(x) = 30 .
\]
Write your answer as a solution set. 

\item Solve the equation 
\[
       f(x) = 31 .
\]
Write your answer as a solution set. 

\item Solve the equation
\[
      f(x) = f(10).
\]
Write your answer as a solution set.  


\end{enumerate}
\end{question}


\section{Completing the Square}

\begin{question} \label{Qldf3fr3fd}

\begin{enumerate}
\item Distribute: $(x+5)^2$

\item Distribute: $(x-6)^2$
\end{enumerate}
\end{question}

\begin{question} \label{Q9r3mmzzzz}
Let
\[
       g(x) = x^2 + 8x + 7 .
\]
\begin{enumerate}
\item Solve the equation
\[
    g(x) = 0 .
\]
Write your answer as a solution set.

\item Solve the equation
\[
    g(x) = 7 .
\]
Write your answer as a solution set.

\item Solve the equation
\[
    g(x) = 91 
\]
by completing the square. Write your answer as a solution set.

\end{enumerate}
\end{question}


\section{Twice as Far}

\begin{question} \label{Qdr1000}

You'll learn about tangent lines to a curve in calculus, but for now we'll say that a line is tangent to a circle if it intersects the circle in exactly one point.


\begin{onlineOnly}
    \begin{center}
\desmos{aosglrotza}{900}{600}
\end{center}
\end{onlineOnly}

\href{https://www.desmos.com/calculator/aosglrotza}{141: Circles tangent to lines}

\end{question}

\end{document}
